%% ============================================================
%%  Chapter 5 — Discussion
%% ============================================================
\chapter{Discussion}
\label{ch:discussion}

This chapter interprets the experimental results in the context of the
research objectives, analyses the interpretability--performance trade-off,
discusses limitations of the simulation framework, and examines key
design decisions.
Section~\ref{sec:disc_sedm} analyses SEDM behaviour and the role of the
Markov chain, extended with an interpretation of the traffic-realism and
escalation-tracking results.
Section~\ref{sec:disc_dqn} interprets DQN training dynamics.
Section~\ref{sec:disc_tradeoff} addresses the interpretability--performance
trade-off and directly revisits whether a learning-based mechanism should
have been used for dynamic response.
Section~\ref{sec:disc_limitations} discusses limitations, including those
specific to the new mechanisms.
Section~\ref{sec:disc_design} examines design choices.
Section~\ref{sec:future} outlines future work directions.

%% ── 5.1 SEDM Analysis ────────────────────────────────────────
\section{SEDM Performance and Markov Chain Adaptation}
\label{sec:disc_sedm}

\subsection{Why $>$\,99\% Detection is Achievable}

The near-perfect detection rates across all four intent profiles are
remarkable given that the SEDM has no trainable parameters and requires
no labelled data.
Three structural factors explain this result.

\textbf{Factor 1: Evaluation protocol.}
The 200-step evaluation episodes involve continuous attack traffic; there
are no extended benign periods.
Because the kill chain stages visited in all four intents are predominantly
Installation, C2, and Actions on Objectives (as shown by the kill chain
distribution in Chapter 4, \S4.4), the SEDM's matrix almost always returns
BLOCK or ALERT --- both of which count as detections.
In other words, the high detection rate is partly a consequence of the
evaluation scenario: it is easy to achieve high detection when attacks are
dense and the kill chain has already progressed.

\textbf{Factor 2: Conservative matrix design.}
The SEDM never returns ALLOW for stages above RECONNAISSANCE under any
escalation risk band.
Even in the Low band at WEAPONIZATION, the action is LOG rather than ALLOW.
This conservatism ensures that any step beyond RECONNAISSANCE is responded
to with at least a LOG action, which counts as a detection.

\textbf{Factor 3: Intent-adaptive escalation risk.}
The Markov chain escalation risk computation (\S3.4, Step 1) correctly
classifies each intent's typical trajectory into the appropriate band.
For STEALTHY at RECONNAISSANCE, the low escalation risk correctly returns
ALLOW (or LOG in the Medium band), matching the low-threat nature of the
traffic.
For AGGRESSIVE at EXPLOITATION, the high escalation risk correctly returns
BLOCK or ALERT, matching the high-threat nature of the traffic.
The alignment between the Markov model and the evaluation scenarios means
that the band classification is correct for the vast majority of steps.

\subsection{False Positive Rate: Structural Causes}

Under the corrected evaluation alignment (each action scored against the
ground-truth label of the observation it was actually chosen from, rather
than the following step's label), the measured false positive rate is
0.00\% across all four intents under classifier-driven decisions. An
earlier draft of this evaluation, produced before the one-step label-lag
bug (documented as Bug 8/9 in \texttt{docs/BUGS\_AND\_FIXES.md}) was fixed,
reported false positive rates varying by a factor of $\sim$10 across
intents (3.33\% for TARGETED vs.\ 35.56\% for STEALTHY); that variation did
not reflect real SEDM behaviour and has been withdrawn.

The 0.00\% figure is not a claim that the SEDM is immune to false positives
in general.
It follows from two specific, checked properties of the current
experimental setup: (1) the R1 override unconditionally maps a
predicted-NORMAL sample to ALLOW, so any false positive must originate
upstream in the classifier mistaking a true NORMAL sample for an attack
type; and (2) the Random Forest classifier achieves 100\% precision and
recall on the NORMAL class specifically (99.85\% accuracy overall), because
the synthetic feature simulator draws each attack type from a disjoint
parametric distribution by construction.
Because real network telemetry does not offer this clean separability, a
feature-noise robustness sweep (Chapter 4, \S4.8) injects multiplicative
Gaussian noise onto the classifier's input features and recovers a
non-zero, monotonically increasing false-positive rate (0.00\% at
$\le$10\% noise, rising to 10.00\% at 50\% noise).
This sweep, not the noiseless 0.00\% headline figure in isolation, is the
evidence that should inform expectations about false-positive behaviour
against real, noisy traffic.

\textbf{A statistical, not just a mechanistic, qualification.}
The two structural properties above explain \emph{why} the classifier-driven
pipeline produced zero observed false positives, but they leave open a
separate question the original evaluation did not directly address: how
much evidence does a 0.00\% figure computed from the original protocol's
traffic actually provide?
\S4.1's revised primary protocol (Table 4.14) answers this directly by
counting the benign-labelled samples the 0.00\% figure was computed from
--- 33 to 114 per intent, since \texttt{benign\_ratio}\,=\,0 injects no
benign traffic and these samples are only whatever the attacker's own
Markov chain happened to produce.
A Wilson 95\% confidence interval over so few samples has an upper bound as
high as 10.43 percentage points (TARGETED); the 0.00\% point estimate was
therefore statistically compatible with a true false-positive rate an order
of magnitude higher than the headline number suggests, simply because too
few benign samples existed to observe one.
Re-running the identical decision logic with \texttt{benign\_ratio}\,=\,0.3
(50 episodes $\times$ 500 steps, the same budget already justified in
Table 4.8) multiplies the benign sample count by 70--230$\times$ and
resolves this ambiguity: the measured false-positive rate is a small but
genuinely non-zero 0.60--0.79\%, with confidence intervals narrow enough
(roughly 0.2--0.4 percentage points) to be informative, and independently
corroborated by the specificity figures in Table 4.15 (0.992--0.995)
computed via unrelated evaluation infrastructure.
This does not overturn Table 4.1 --- detection rate is essentially
unchanged between protocols, and Table 4.1 remains a valid characterisation
of SEDM behaviour under continuous-attack traffic --- but it does mean the
0.00\% headline figure should never have been read as evidence of
near-immunity to false positives; it was, in retrospect, primarily a
statement about how rarely the original protocol exercised the benign case
at all.

\subsection{Reward Ordering: TARGETED $>$ AGGRESSIVE $>$ STEALTHY $>$ OPPORTUNISTIC}

The reward ordering across intents (Chapter 4, Table 4.1) reflects the
alignment between the SEDM's action preferences and the reward matrix.

TARGETED receives the highest reward because its concentrated exploitation
path consistently generates high-threat steps that are met with ALERT
(reward $+6.0$ for critical traffic) with very few false positives (ALLOW
on benign, $+1.0$) or suboptimal responses (TROLL on critical, $+0.5$).

OPPORTUNISTIC receives the lowest reward because its scattered attack types
frequently produce medium-severity steps that receive TROLL ($+3.0$ for
medium) or BLOCK ($+1.5$ for medium) when ALERT would yield higher reward
for critical steps.
The SEDM does not see the future of the OPPORTUNISTIC campaign and
calibrates to the current stage's escalation risk, which for medium-band
Delivery-stage steps recommends TROLL.
If the same step were at INSTALLATION with high escalation risk, ALERT
would be recommended and the reward would be higher.

\subsection{Determinism vs.\ Stochasticity}

The SEDM is fully deterministic: given the same state (and, in this
version, the same reputation input), it always returns the same action.
This determinism is a feature for operational deployability (the policy is
predictable and auditable) but could be a vulnerability against an adaptive
attacker who can fingerprint the policy and craft states that trigger
suboptimal responses.

The override rules (particularly R3, which responds to high attack
frequency, and --- new in this version --- R4, which responds to
accumulated cross-session reputation) partially mitigate this by creating
response modes that adapt to the tempo and history of the campaign, even
if the per-step action for a \emph{given} state remains fixed.
\S4.10 shows, however, that R3's practical firing behaviour is itself
highly sensitive to which escalation signal drives it, and an adversary who
maintains escalation rate just below the active threshold would still
avoid triggering R3 under either signal.
R4 changes this calculus for a returning adversary specifically: \S4.11
shows that after a small, deterministic number of prior offenses, R4 fires
regardless of how the \emph{current} event is crafted, closing off the
specific evasion strategy of ``look identical to the current event on
every visit'' for a source with sufficient accumulated history --- though
not, as \S3.10.1 and Section~\ref{sec:disc_limitations} discuss, the
strategy of simply never returning as the same identifiable source.
Future work should investigate whether small perturbations to the
threshold values or the addition of a randomised component would further
improve robustness against adaptive adversaries.

\subsection{Interpreting the Traffic-Realism Results (\S4.9)}
\label{sec:disc_realism}

The 0.45 percentage-point classifier accuracy drop measured in \S4.9 is
best read as a validity check succeeding rather than a regression.
Before this extension, both the classifier's 99.85\% accuracy and the
SEDM's 0.00\% headline false-positive rate depended on a specific,
load-bearing assumption: that every sample of a given class is drawn
independently from the same fixed distribution as every other sample of
that class, with no session-level or persona-level structure.
Real network sources do not satisfy that assumption --- a single session's
traffic is autocorrelated with itself, and ``benign'' traffic is not one
behavioural population but several.
\S4.9's result confirms that when this assumption is relaxed even
modestly, in a direction chosen for realism rather than adversarial
difficulty, accuracy moves in the expected direction (down) and by an
amount consistent with the specific classes affected (EXPLOITS/GENERIC, and
NORMAL specifically).
This is the same qualitative conclusion the feature-noise robustness sweep
(\S4.8) already established via artificial noise injection, now obtained
from a structural change to the generator with no noise injected at all
--- two independent probes of the same underlying claim (near-zero false
positives is an artefact of separability, not a property of the SEDM
policy) agreeing with each other is stronger evidence than either alone.

\subsection{Interpreting the Escalation-Tracking Results (\S4.10)}
\label{sec:disc_escalation_mode}

The R3 trigger-rate gap between window and EMA tracking
(44.7--90.0\% vs.\ 0.7--22.0\%, \S4.10) is the clearest demonstration in
this thesis that a design decision motivated by principled reasoning
(\S3.9's argument that severity-weighting is more informative than
frequency-counting) can also have large, non-obvious operational
consequences that only become visible empirically.
The window signal's near-universal saturation under continuous attack
traffic means that, under the original design, R3 functions less as a rare
escalation trigger and more as a background condition that happens to be
true for most of a typical episode --- a materially different operating
characteristic from the one implied by \texttt{RATE\_THRESHOLD}'s design
rationale in \S4.8 (``fires only under sustained, dense attack traffic'').
Neither this thesis's original evaluation (\S4.1--\S4.8) nor its underlying
design intent were wrong about \emph{what} R3 is supposed to do; \S4.10
reveals that the window signal, specifically under the continuous-attack
protocol used throughout most of this thesis's evaluation, does not
achieve that intent as cleanly as the design rationale assumed.
This finding directly motivates re-examining \texttt{RATE\_THRESHOLD}'s
calibration (Section~\ref{sec:future}) and explains the bounded threshold
controller's regime-dependent behaviour reported in \S4.12.

%% ── 5.2 DQN Training Dynamics ────────────────────────────────
\section{DQN Training Dynamics Interpretation}
\label{sec:disc_dqn}

\subsection{Rapid Detection Rate Improvement}

The detection rate jump from 87.6\% (episode 0) to 97.4\% (episode 1)
reflects the dense attack structure of the OPPORTUNISTIC training scenario
combined with the DQN's rapid initial learning.
After episode 0, the replay buffer contains 500 transitions covering a wide
range of states; the first gradient update trains the network to associate
non-ALLOW actions with positive rewards for attack states.
This coarse association suffices to achieve $>$\,97\% detection because the
vast majority of states in OPPORTUNISTIC campaigns are high-threat attack
states where any non-ALLOW action yields positive reward.

The implication is that the apparent high detection rate of the DQN does not
reflect a nuanced understanding of the threat landscape; rather, it reflects
a learned bias toward non-ALLOW responses for the dense attack distribution
of OPPORTUNISTIC traffic.
This interpretation is supported by the near-constant false-positive rate
of 1.0 throughout training: the DQN has effectively learned to never ALLOW,
which maximises reward on the training distribution but would fail on a more
balanced traffic mixture.
This specific failure mode --- a policy whose apparent competence is a
distributional artefact rather than genuine discrimination --- is the
central piece of project-specific evidence Section~\ref{sec:disc_tradeoff}
relies on when arguing against a learning-based mechanism for the
dynamic-response extensions introduced in this version.

\subsection{Loss Dynamics}

The increasing loss trend through training (from $\approx$\,0.97 to
$\approx$\,2.4) is not indicative of divergence.
It reflects the expansion of the Q-value range as the policy learns to
differentiate high-reward from low-reward state-action pairs.
The Bellman target $y_i = r_i + \gamma \max_{a'} Q(s', a'; \theta^-)$
grows as the Q-function improves, causing the absolute Huber loss to
increase even as the relative TD error decreases.
The stable variance in the later training phase (no upward spikes after
episode 100) confirms that the stabilising mechanisms (target network,
gradient clipping) are effective.

\subsection{DQN Limitations in This Setting}

The DQN baseline reveals two limitations specific to the HoneyIQ setting.

\textbf{Class imbalance in training data.}
The OPPORTUNISTIC training episodes contain predominantly attack traffic.
The DQN is never exposed to episodes with significant benign traffic,
so it cannot learn to discriminate benign from malicious states.
The effective strategy of ``always respond with a non-ALLOW action'' achieves
high training reward but is not a robust policy.

\textbf{Transfer gap.}
Training exclusively on OPPORTUNISTIC intent means the Q-function is
calibrated to the specific attack distribution of OPPORTUNISTIC campaigns.
Evaluating against STEALTHY or TARGETED would require the DQN to generalise
to different threat-level time series; without intent-conditioned training
or multi-intent training, this generalisation is likely to be imperfect.
The SEDM avoids this limitation entirely through the analytical escalation
risk computation.

Both limitations are properties of the \emph{training distribution design},
not inherent to DQN or deep RL as a class of methods --- a better-designed
training protocol (benign-heavy episodes, multi-intent training) would very
plausibly avoid them.
This distinction matters for Section~\ref{sec:disc_tradeoff}: the argument
developed there is not that DQN is categorically inferior, but that
\emph{this specific, already-executed} training run demonstrates a concrete
failure mode this project would need to guard against before trusting a
learned mechanism for anything, and that guarding against it (redesigning
the training distribution, and --- separately --- establishing a
trustworthy reward signal, discussed below) is itself substantial,
currently unstarted work.

%% ── 5.3 Interpretability–Performance Trade-off ───────────────
\section{Interpretability and Performance Trade-off}
\label{sec:disc_tradeoff}

\subsection{The Case for SEDM}

The SEDM achieves $>$\,99\% detection across all four intent profiles with
zero training overhead, full transparency, and an auditable decision
procedure.
These properties are highly valuable for operational deployment:
%
\begin{itemize}
  \item \textbf{No training data requirement.}
        The SEDM requires only the Markov transition model, which encodes
        assumed attacker behaviour.
        It can be deployed immediately without a historical attack corpus.
  \item \textbf{Analytical escalation risk.}
        The escalation risk is computed from the transition matrix rather
        than estimated from data, making it exact under the model assumptions
        and independent of observation noise.
  \item \textbf{Full traceability.}
        Every action can be traced to a specific matrix cell and override
        rule.
        An analyst can reproduce any decision with pencil and paper ---
        including, in this version, an R4-driven decision, which reduces to
        one decayed score and one threshold comparison (\S3.10.1, \S4.11).
  \item \textbf{Easy modification.}
        Security practitioners can update the matrix entries to reflect new
        threat intelligence (e.g., raising the response at WEAPONIZATION
        for a specific campaign) without retraining any model.
\end{itemize}

\subsection{The Case for DQN}

The DQN baseline offers adaptive flexibility that the SEDM cannot provide:
%
\begin{itemize}
  \item \textbf{Reward-optimised responses.}
        Given sufficient training data and a well-designed reward function,
        the DQN can discover optimal action strategies that are not
        immediately obvious to human designers.
  \item \textbf{Continuous adaptation.}
        Online DQN training (with a sliding replay buffer) allows the policy
        to adapt as the attack distribution evolves, provided new transitions
        are added to the buffer.
  \item \textbf{Nuanced feature sensitivity.}
        In a deployment with richer state representations (e.g., including
        classifier probability vectors, connection metadata), the DQN can
        exploit fine-grained patterns that a fixed matrix cannot capture.
\end{itemize}

\subsection{Revisiting DQN for Dynamic Response}
\label{sec:disc_dqn_revisit}

This version of the system introduces exactly the kind of runtime
adaptivity the ``case for DQN'' above describes as DQN's comparative
advantage: behaviour that changes in response to accumulated experience
(\S3.10.1) and to observed operating conditions (\S3.10.2).
It is therefore the natural point to ask directly: given that this project
already has a trained DQN and its full training record, should \emph{that}
mechanism --- or a similar learned one --- have been used for \S3.10's
dynamic-response extensions instead of the auditable mechanisms actually
built?
The answer developed here is no, for reasons specific to this project's own
evidence rather than a general anti-RL position.

\textbf{The DQN's own recorded failure mode is directly relevant.}
Section~\ref{sec:disc_dqn} already establishes that the trained DQN's
apparent competence --- detection rate rising from 87.6\% to 97.4\% within a
single episode --- was a distributional artefact: the agent learned
``almost never ALLOW'' because its training distribution made that strategy
reward-maximising, not because it learned genuine threat discrimination.
A learning-based dynamic-adjustment mechanism trained on this project's
synthetic traffic would face the identical risk: absent a deliberately
corrected training distribution (multi-intent, benign-heavy), there is
every reason grounded in this project's own results to expect it would
converge to an analogous shortcut rather than the nuanced adjustment
behaviour actually wanted.

\textbf{The evaluation harness's own defect history compounds this risk.}
\texttt{docs/BUGS\_AND\_FIXES.md} documents two independent instances (Bugs
8 and 9) of a one-step label/observation misalignment that silently
inflated measured false-positive rates by up to an order of magnitude
before being caught --- in code that only needed to compute a descriptive
statistic \emph{after} an episode had concluded.
A reward signal for a learning agent would be computed by structurally
similar logic, on every single step, directly shaping what the agent learns
rather than merely mis-reporting a number afterward.
Given that a comparatively simple defect class evaded review twice in this
codebase, extending trust to a reward-shaped training loop --- a strictly
higher-stakes use of the same kind of step-alignment logic --- is not
warranted without a materially more adversarial validation pass on the
reward computation than has been performed here.

\textbf{No trustworthy reward signal exists for the specific mechanisms in
question.}
\S3.10.2 already establishes, as a design constraint rather than an
afterthought, that no ground-truth feedback pathway exists in this system
for whether an individual R3 trigger was operationally correct; the same is
true for R4.
A DQN (or any RL method) requires exactly such a signal.
Manufacturing one from a proxy --- R1 or R3's own trigger frequency, for
instance --- would not constitute genuine ground truth, and would risk
reproducing the DQN's documented failure mode in a new location: a policy
that optimises a measurable proxy while providing no evidence of solving
the actual underlying problem.

\textbf{This reasoning extends the recommendation already reached in the
original version of this thesis.}
Citing \citet{rudin2019stop}: prefer an interpretable model until a
black-box alternative demonstrates clear, validated superiority \emph{and}
can itself be explained after the fact.
The two mechanisms actually built (\S3.10) satisfy that standard directly
and by construction --- both reduce, at any instant, to a small number of
named numeric quantities (\S3.10.3) --- rather than requiring a post-hoc
explanation layer bolted onto an opaque function approximator.

\textbf{What would change this conclusion.}
If a learning component is judged worth pursuing in future work, the most
defensible form identified is a small contextual bandit --- not a
Q-network, and not a general-purpose RL agent --- selecting among a handful
of preset \texttt{RATE\_THRESHOLD} values, with a reward computed from a
genuinely trustworthy signal: an analyst's own confirmed true/false-positive
label on a specific logged decision (available in principle from
\texttt{DummyHoneypot}'s action audit log, \S3.10.1).
That labelling pathway --- closing the loop from a human judgement back
into a trainable signal --- does not currently exist anywhere in this
system.
Section~\ref{sec:future} identifies building it as the specific, concrete
prerequisite that would need to be satisfied before this section's
conclusion should be revisited, rather than treating ``more data'' or
``more training'' in the abstract as sufficient.

\subsection{Practical Recommendation}

The results suggest a two-stage deployment strategy, now with the specific
mechanisms available to implement the second stage without abandoning the
first:
%
\begin{enumerate}
  \item \textbf{Initial deployment}: Use the SEDM, including the R4
        reputation override and (optionally) the bounded R3 threshold
        controller, as the primary policy.
        Its immediate operability, interpretability, and strong baseline
        performance ($>$\,99\% detection) make it the appropriate default,
        and both dynamic-response extensions preserve this property by
        construction.
  \item \textbf{Refinement, gated on a real feedback pathway}: Train a DQN
        agent alongside the SEDM using live honeypot interactions only once
        a genuine analyst-labelling pathway exists.
        If the DQN demonstrably outperforms the SEDM on detection rate or
        false positive rate under rigorous evaluation, and its decisions
        can be explained post-hoc via SHAP~\citep{lundberg2017unified} or
        similar techniques, it can be used to augment or eventually replace
        specific components of the SEDM for specific intent profiles.
\end{enumerate}

This strategy mirrors the human-AI collaboration model recommended by
\citet{rudin2019stop} for high-stakes decision-making: start with an
interpretable model, and only adopt a black-box model when there is
clear evidence of superior performance and adequate explanation mechanisms.
The addition in this version is a specific, actionable gate on stage 2 ---
the analyst-labelling pathway --- rather than leaving ``sufficient training
data'' as an unspecified future condition.

%% ── 5.4 Limitations ──────────────────────────────────────────
\section{Limitations}
\label{sec:disc_limitations}

\paragraph{Synthetic data only.}
Feature distributions are parametric approximations designed to reproduce
the qualitative signatures of UNSW-NB15 attack types.
They do not capture long-range temporal correlations, protocol interactions,
or the noise characteristics of real network traffic.
The simulation-to-reality gap means that results cannot be directly
extrapolated to live network deployments without validation on real
honeypot logs.
The traffic-realism extensions in \S3.2.4 narrow this gap slightly (session
coherence, persona diversity) but do not close it; they remain parametric
approximations, not fits to real data.

\paragraph{Defender action does not affect attacker trajectory.}
In the current formulation, the attacker's Markov chain evolves
independently of the defender's actions.
In practice, BLOCK or ALERT responses may deter, delay, or redirect an
attacker.
Modelling the attacker as a best-responding agent --- using multi-agent RL,
game-theoretic formulations, or a secondary attacker Markov chain that
conditions on the observed defender actions --- would produce more realistic
adversarial dynamics.

\paragraph{Classifier integration is validated only on synthetic data.}
The primary and extended evaluations in Chapter 4 now run the SEDM in a
classifier-driven mode, deciding from the Random Forest classifier's
\emph{predicted} attack type while scoring against ground truth, alongside
an oracle mode that decides from ground-truth labels directly.
This confirms that the SEDM remains robust (detection $\ge$\,99.9\%,
classification metrics $\ge$\,0.999) under the classifier's residual
misclassification rate.
What remains untested is robustness under \emph{realistic} classification
noise: the classifier's near-perfect accuracy (99.85\% originally, 99.40\%
under the traffic-realism extensions of \S4.9) is itself a property of the
synthetic feature simulator's clean class separability.
The feature-noise robustness sweep (\S4.8) and the traffic-realism result
(\S4.9) both show this ceiling degrades predictably as input fidelity or
class separability decreases, but the SEDM's escalation-risk computation
has not been re-evaluated end-to-end under noisy classifier input --- only
the classifier's raw accuracy has been.
Closing this gap, and eventually validating against a classifier trained on
real (non-synthetic) network features, is the genuinely open item here.

\paragraph{No network topology.}
HoneyIQ models all interactions as single-session point contacts without
spatial structure.
Real network intrusions involve lateral movement, multi-hop attacks, and
topological constraints that are not captured.
Extending the simulation to a graph-based network topology would increase
realism at the cost of significantly higher state-space complexity.

\paragraph{Static Markov transition matrices.}
The attacker's transition matrices are fixed for each intent and do not
adapt during an episode.
Real adversaries update their strategies based on the defender's responses.
An adversary that observes consistent BLOCK actions on EXPLOITS might switch
to a different attack type or modify its escalation speed.
A dynamic Markov model whose transition matrices evolve in response to
defender feedback would capture this adaptive behaviour.
This is also, precisely, why \S3.10.2 excludes
\texttt{ESC\_LOW\_THRESHOLD}/\texttt{ESC\_HIGH\_THRESHOLD} from the adaptive
mechanisms introduced in this version: those thresholds bucket a quantity
($\rho(k, \pi)$) computed from exactly these static matrices, so no runtime
signal exists that could honestly justify adapting them without first
addressing this limitation.

\paragraph{Episode termination by truncation.}
Episodes terminate only when \texttt{max\_steps} is reached; there is no
natural termination condition corresponding to attacker success or defeat.
In a real scenario, an attacker who achieves Actions on Objectives (data
exfiltration, system destruction) would terminate the campaign.
Modelling goal-conditioned episode termination would allow the evaluation
to capture the probability of attacker success, a more operationally
relevant metric than episode reward.

\paragraph{Reputation constants are unvalidated defaults.}
The \texttt{ReputationTracker} parameters (\texttt{offense\_increment} =
0.25, \texttt{decay\_half\_life} = 6 hours, \texttt{REPUTATION\_THRESHOLD} =
0.60) were chosen to produce a reasonable-looking \emph{shape} of behaviour
--- gradual accumulation, meaningful but non-permanent memory, a threshold
crossing after a small number of offenses (\S4.11) --- rather than fit to
any real attack-frequency or session-return-interval data.
\S4.11 characterises their behaviour precisely, but precision is not the
same as validation: a real deployment would need these constants
recalibrated against the actual return-visit patterns and TTL
characteristics of its own traffic before the specific ``four offenses''
finding in \S4.11 could be relied upon as an operational figure.

\paragraph{The R4-before-R1 ordering is a deliberate but consequential
security-posture choice.}
\S3.10.1 and \S3.4 place R4 ahead of R1 so that a source with sufficient
accumulated reputation cannot regain unconditional trust via a single
benign-looking event.
The direct cost of this choice --- acknowledged there and restated here as
a limitation rather than only a design note --- is that a shared or
dynamically-reassigned address (NAT, DHCP churn, a compromised-then-cleaned
host) that was briefly malicious remains penalised for the reputation
half-life even after genuinely malicious activity from that address has
ceased.
This thesis does not evaluate the frequency with which such address
reassignment would occur in a real deployment, nor therefore the
real-world magnitude of this cost; doing so is identified as future work
(Section~\ref{sec:future}).

\paragraph{The bounded threshold controller's bound is calibrated for
non-saturated conditions.}
\S4.12 shows that \texttt{AdaptiveThresholds}' $\pm 0.10$ bound is
materially more effective under mixed-traffic conditions (7.60\%
$\rightarrow$ 4.40\%) than under the continuous-attack conditions used
throughout most of this thesis's primary evaluation protocol (45.70\%
$\rightarrow$ 44.40\%), because the underlying window-mode escalation
signal is saturated near its ceiling for most of a continuous-attack
episode (\S4.10).
This is not a defect in the controller's logic --- it behaves exactly as
specified in both regimes --- but it does mean the controller's practical
utility, under the default window-mode escalation signal and the default
bound, is regime-dependent in a way that was not obvious prior to
\S4.12's direct measurement.

%% ── 5.5 Design Choices ───────────────────────────────────────
\section{Design Choices and Trade-offs}
\label{sec:disc_design}

\paragraph{LayerNorm over BatchNorm.}
LayerNorm~\citep{ba2016layer} normalises activations across features for
each sample independently, making it batch-size-agnostic.
For DQN training with mini-batches of 64 samples drawn from a diverse
replay buffer, BatchNorm statistics would be computed from a
heterogeneous mix of state types, potentially introducing instability.
LayerNorm avoids this dependency.
Empirically, no training instability attributable to normalisation was
observed during the 300-episode training runs.

\paragraph{Hard target network update.}
A hard copy of the policy network parameters every 150 steps was chosen
over a Polyak soft update ($\theta^- \leftarrow \tau\theta + (1-\tau)\theta^-$,
$\tau \approx 0.005$) for implementation simplicity.
Both approaches stabilise training; the hard update introduces discrete
target shifts every 150 steps but avoids the hyperparameter sensitivity
of choosing $\tau$.
At the episode length of 500 steps, 150-step updates correspond to roughly
three target updates per episode, providing stable targets over a substantial
fraction of each episode.

\paragraph{Composite threat-level weight assignment.}
The weight assignment (attack severity 45\%, kill chain stage 35\%,
escalation rate 15\%, cumulative count 5\%) was designed by domain
reasoning rather than optimisation.
Attack severity receives the highest weight because it is the most direct
indicator of the attacker's current capability and impact.
Kill chain stage receives the second-highest weight because it captures the
campaign context.
The 5\% weight on cumulative attack count is a minor adjustment that
differentiates isolated probes from sustained campaigns without dominating
the threat signal.

An alternative approach would optimise these weights jointly with the reward
matrix using Bayesian optimisation or multi-objective optimisation, with
detection rate and false-positive rate as objectives.
This is left for future work.

\paragraph{SEDM band thresholds.}
The Low/Medium ($\rho = 0.35$) and Medium/High ($\rho = 0.65$) thresholds
were chosen to create roughly equal-width bands in the $[0, 1]$ risk space.
Intent-specific analysis confirms that the thresholds correctly separate
the intent profiles: STEALTHY typically falls in the Low band, AGGRESSIVE
in the High band, and TARGETED/OPPORTUNISTIC in the Medium band for most
stages.
Shifting the thresholds could alter the balance between false positives
(lower thresholds) and false negatives (higher thresholds) for specific
stages.
As established in \S3.10.2 and confirmed empirically throughout Chapter 4,
these two thresholds are the ones deliberately excluded from this
version's adaptive mechanisms, precisely because they have no live
counterpart to adapt against.

\paragraph{Synthetic classifier training data.}
Generating classifier training data from the parametric feature distributions
avoids the need for a real labelled dataset and enables instant re-training
with any desired class balance.
The risk is classifier over-fit to synthetic data that may not generalise
to real traffic.
In the HoneyIQ evaluation, this risk is moot because the evaluation features
are also synthetic; in a real deployment, the classifier would need to be
retrained on real network data.
The independent-sample override introduced in \S3.2.4 (drawing fresh
intensity/persona per training sample rather than reusing one session's
profile) is a design choice specifically aimed at not making this
synthetic-data risk \emph{worse} than necessary: without it, the
classifier's training data would have exhibited an additional, avoidable
form of artificial regularity.

\paragraph{Why R4 precedes R1 rather than the reverse.}
An alternative design would check R1 first, letting any benign-looking
event through regardless of reputation, and reserve R4 for events that R1,
R2, and R3 all decline to escalate.
This alternative was considered and rejected: it would make R4 functionally
inert for exactly the scenario it exists to address (a returning,
previously-malicious source opening with an innocuous probe), since R1
would already have resolved the decision to ALLOW before R4 is ever
consulted.
Checking R4 first is the only ordering under which the mechanism can do the
job stated in its own design rationale (\S3.10.1); the security-posture
cost of that ordering is acknowledged directly in Section~\ref{sec:disc_limitations}
rather than treated as a free choice.

\paragraph{Why the threshold controller is scoped to \texttt{RATE\_THRESHOLD}
alone.}
\S3.10.2 explains the reasoning; the design-choice framing is that this was
a decision to under-promise rather than over-build.
A more ambitious design might have attempted to adapt the SEDM table
entries themselves, or the band thresholds, in response to some composite
signal.
Both were rejected specifically because no genuine, non-circular signal
exists in this system to adapt them against --- building such a mechanism
would have produced something that \emph{looks} more sophisticated while
being less trustworthy than the narrowly-scoped controller actually built,
a trade-off this thesis consistently resolves in favour of narrower scope
and stronger auditability (\S2.8, Section~\ref{sec:disc_tradeoff}).

%% ── 5.6 Future Work ──────────────────────────────────────────
\section{Future Work Directions}
\label{sec:future}

\subsection{Classifier Integration Under Realistic Noise}
\label{sec:future_classifier}

The SEDM's decision pipeline already supports deciding from the Random
Forest classifier's predicted attack type rather than ground truth
(Chapter 4), and detection remains $>$\,99.9\% under this classifier-driven
mode.
The most immediate remaining extension is to re-run this evaluation with
the classifier's input features corrupted by the same realistic
measurement noise used in the feature-noise robustness sweep (\S4.8), to
test whether the SEDM's escalation-risk computation and downstream action
distribution degrade gracefully as classifier accuracy falls from 99.85\%
(or 99.40\% under the \S3.2.4 extensions) toward the 87.65\% observed at
50\% feature noise.
A confidence-weighted band assignment --- where uncertain classifier
predictions widen the escalation risk estimate --- could provide a principled
way to handle the resulting ambiguous observations.

\subsection{Adaptive Attacker Modelling}

Modelling the attacker as a best-responding agent (e.g., using multi-agent
RL or a game-theoretic Stackelberg formulation) would introduce genuine
adversarial dynamics.
An attacker trained against the SEDM might learn to manipulate the
escalation risk computation by dwelling at low-risk stages longer than
the Markov chain predicts, or by generating traffic that bypasses the
override rules --- including, in this version, an attacker that
deliberately avoids ever reusing a source identity, which would defeat R4
(\S3.10.1) by construction.
Evaluating SEDM robustness against such an adaptive adversary would provide
a stronger performance guarantee.

\subsection{Multi-Intent Training for DQN}

Exposing the DQN to a mixture of intent profiles during training --- sampling
a random intent at the start of each episode --- would broaden the state
distribution and reduce the tendency to over-specialise on OPPORTUNISTIC
traffic.
An intent-conditioned DQN, where the current intent is provided as input
(as it is in the 24-dimensional state vector), could be compared with the
SEDM on cross-intent generalisation.
This remains a prerequisite, not a completed step, for the kind of
learning-based dynamic-response mechanism discussed and set aside in
Section~\ref{sec:disc_tradeoff}.

\subsection{Real Network Feature Distributions}

Fitting the 15 feature distributions to the actual UNSW-NB15 records
(using maximum likelihood estimation or kernel density estimation per
attack category) would reduce the simulation-to-reality gap further than
the session-coherence and persona extensions of \S3.2.4 do.
Alternatively, a conditional generative model (e.g., a conditional VAE
or normalising flow) trained on real packet data would enable richer,
correlated feature generation.

\subsection{Network Topology Extension}

Extending HoneyIQ to a multi-node network topology with explicit honeypot
placement decisions would introduce a spatial dimension to the problem.
The SEDM could be extended to consider both the kill chain stage and the
network position of the attacker, enabling topology-aware response policies.

\subsection{Distributional Reinforcement Learning}

Replacing the standard DQN with a distributional RL variant
(e.g., C51~\citep{wang2016dueling} or QR-DQN) would allow the agent to
maintain a distribution over returns rather than a point estimate.
The resulting uncertainty estimates could improve decisions at threat-band
boundaries, where the optimal action is least clear.
As with multi-intent training, this remains gated on the
training-distribution and reward-trustworthiness prerequisites identified
in Section~\ref{sec:disc_tradeoff} before it should be considered for
anything beyond offline experimentation.

\subsection{OpenCanary Integration}

HoneyIQ includes a live OpenCanary-integration pipeline (\S3.6, \S3.10)
that shares its state representation and decision policy with the
synthetic training environment.
Replaying real OpenCanary session logs through this pipeline --- rather
than only the synthetic emulator currently used to exercise it (\S4.11) ---
would provide an out-of-distribution test of generalisation to genuine
attack traffic, bridging the simulation-to-reality gap without requiring
live deployment, and would additionally provide the first opportunity to
validate the \texttt{ReputationTracker} constants
(Section~\ref{sec:disc_limitations}) against real source-return-interval
statistics.

\subsection{An Analyst-Labelling Pathway}

Section~\ref{sec:disc_tradeoff} identifies this as the single concrete
prerequisite for any future learning-based dynamic-adjustment mechanism: a
pathway by which a human analyst's confirmed true/false-positive judgement
on a specific logged decision (available from the existing action audit
log, \S3.10.1) is captured and made available as a training signal.
Building this pathway is itself independent of any specific learning
algorithm choice, and should precede --- not accompany --- any decision to
revisit the conclusion reached in Section~\ref{sec:disc_tradeoff}.

\subsection{Re-Calibrating \texttt{RATE\_THRESHOLD} for EMA Tracking and
Mixed Traffic}

\S4.10 and \S4.12 together show that \texttt{RATE\_THRESHOLD}'s current
value (0.80) and the bounded controller's default bound ($\pm 0.10$) were
both, implicitly, calibrated against window-mode escalation tracking under
continuous-attack conditions --- the protocol used throughout most of this
thesis's evaluation --- and behave differently under EMA tracking or
mixed-traffic conditions.
An independent recalibration pass, repeating the standard-error-based
methodology already used for episode/step counts (\S4.8, Table 4.8) but
applied to \texttt{RATE\_THRESHOLD} and the controller's
\texttt{target\_rate}/\texttt{bound} under each signal-and-traffic-regime
combination, would place both on the same empirically-justified footing as
the parameters already validated in \S4.8.
